% !TeX root = ../main.tex
\chapter{Inversão de dependência: repositório e caso de uso}
\label{ap:dip}

Este apêndice exemplifica a aplicação do princípio da inversão de dependência (SOLID) na arquitetura da API: o caso de uso depende de abstrações de repositório definidas no domínio, e não de implementações concretas de persistência.

\section{Contrato de repositório}

A Listagem~\ref{lst:donation-repo} apresenta a interface \codigo{DonationRepositoryInterface}, situada na camada \codigo{domain}. Operações de leitura e escrita são expressas em termos do agregado \codigo{Donation} e do objeto de valor \codigo{DomainID}, sem menção a MongoDB ou a \textit{schemas} de persistência.

\lstinputlisting[
  language=Java,
  caption={Interface \texttt{DonationRepositoryInterface}},
  label={lst:donation-repo}
]{apendices/codigo/DonationRepositoryInterface.java}
\fonte{Elaborado pelos próprios autores (2026).}

\section{Orquestração no caso de uso}

A Listagem~\ref{lst:complete-pending} mostra o método \codigo{execute} de \codigo{CompletePendingDonationUseCase}. O fluxo valida a entrada, carrega a doação via repositório, verifica \textit{ownership}, delega a transição de estado ao agregado (\codigo{donation.complete}), atualiza o doador e persiste a doação concluída. O cumprimento de metas não é gravado nesse momento: a leitura seguinte monta o \textit{snapshot} FIFO a partir do banco.

\lstinputlisting[
  language=Java,
  caption={Método \texttt{CompletePendingDonationUseCase.execute}},
  label={lst:complete-pending}
]{apendices/codigo/CompletePendingDonationUseCase.java}
\fonte{Elaborado pelos próprios autores (2026).}
